%

\pol{\section{Historia}}
\ita{\section{Storia}}

\pol{Nie będziemy zanudzać Czytelnika sto sześćdziesięcioma trzema wzorami redukcyjnymi, na funkcje i kofunkcje, na sumy i różnice kątów, blablabla.}
\ita{Non annoieremo il lettore con centosessantatré formule di riduzione, per funzioni e cofunzioni, per somme e differenze di angoli, e così via.}
\pol{Jest pełno innych książek, które to robią.}
\ita{Esistono già molti altri libri che lo fanno.}
\pol{Opowiemy trochę o etymologii różnych słów i tym, jak trygonometria rozwijać się będzie na przestrzeni wieków, co powinno być dużo ciekawsze.}
\ita{Racconteremo invece qualcosa sull'etimologia di vari termini e su come la trigonometria si sia sviluppata nel corso dei secoli, il che dovrebbe essere molto più interessante.}

\pol{Współczesne ,,sinus'' i ,,kosinus'' wywodzą się z błędu w tłumaczeniu!}
\ita{I moderni termini ``seno'' e ``coseno'' derivano da un errore di traduzione!}
\pol{Na początku będzie sanskryt i~\emph{jyā} (cięciwa łuku, ze starogreckiego χορδή -- sznurek, struna), a potem transkrypcja do arabskiego \emph{jība}, gdzie pomija się samogłoski (!) i prowadzi do jb, słowa bez żadnego znaczenia.}
\ita{All'inizio vi è il sanscrito \emph{jyā} (corda di un arco; dal greco antico χορδή, cioè corda o stringa), poi la trascrizione araba \emph{jība}, nella quale si omettono le vocali (!) e si ottiene jb, una parola priva di significato.}
% (جب)
\pol{Jb zinterpretują jako \emph{jayb}, które już coś znaczy (po angielsku: ,,bosom, pocket, fold''); Gerard z~Cremony przetłumaczy to do łacińskiego \emph{sinus} (czyli ,,zatoki'', ale też ,,zwisającego fałdu togi na piersi'')...}
\ita{Jb verrà interpretato come \emph{jayb}, che invece ha un significato (in inglese: ``bosom, pocket, fold''); Gerardo da Cremona lo tradurrà in latino come \emph{sinus} (cioè ``baia'', ma anche ``piega della toga sul petto'')...}
\pol{Kosinus to skrót od \emph{complementi sinus} jak w \emph{Canon triangulorum} Edmunda Guntera z 1620 roku; po drodze będzie też Fibonacci oraz jego \emph{sinus rectus arcus}, co pomoże w przyjęciu się ostatecznej wersji.}
\ita{Coseno è un'abbreviazione di \emph{complementi sinus}, come nel \emph{Canon triangulorum} di Edmund Gunter del 1620; nel mezzo vi saranno anche Fibonacci e il suo \emph{sinus rectus arcus}, che contribuiranno all'affermazione della forma definitiva.}

\pol{Nazwy dwóch kolejnych funkcji trygonometrycznych także mają korzenie w łacinie, bowiem \emph{tangens} znaczy ,,dotykający'', zaś \emph{secans} to ,,przecinający''.}
\ita{Anche i nomi di altre due funzioni trigonometriche hanno radici latine: \emph{tangens} significa ``che tocca'', mentre \emph{secans} significa ``che taglia''.}
\pol{Ma to sens, jeśli popatrzymy na okrąg jednostkowy.}
\ita{Ciò ha senso se si osserva il cerchio unitario.}
\pol{Wreszcie termin ,,trygonometria'' bierze się z greckiego τρίγωνον (trigonon -- trójkąt) oraz μέτρον (metron -- miara).}
\ita{Infine, il termine ``trigonometria'' deriva dal greco τρίγωνον (\emph{trigonon}, triangolo) e μέτρον (\emph{metron}, misura).}
% TODO: https://en.wikipedia.org/wiki/History_of_trigonometry#Ancient_Near_East
% TODO: https://en.wikipedia.org/wiki/History_of_trigonometry#Classical_antiquity
% TODO: https://en.wikipedia.org/wiki/History_of_trigonometry#Indian_mathematics
% TODO: https://en.wikipedia.org/wiki/History_of_trigonometry#Chinese_mathematics
% TODO: https://en.wikipedia.org/wiki/History_of_trigonometry#Islamic_world
% TODO: https://en.wikipedia.org/wiki/History_of_trigonometry#European_renaissance
% TODO: https://en.wikipedia.org/wiki/History_of_trigonometry#Modern

\pol{Podstawy trygonometrii rozwinie Hipparchos z Nikei: około 140 r.p.n.e opracuje tablicę cięciw, czyli odpowiednik współczesnych sinusów.}
\ita{Le basi della trigonometria furono sviluppate da Ipparco di Nicea: intorno al 140 a.C. compilò una tavola delle corde, equivalente agli attuali seni.}
\pol{Pierwsza księga poświęcona temu niezależnie od astronomii, \emph{,,Kitāb al-Shakl al-qattā''} (\emph{Rozprawa o czworobokach}), wyjdzie spod pióra Nasira al Dina al-Tusiego około 1250 roku.}
\ita{Il primo libro dedicato all'argomento indipendentemente dall'astronomia, \emph{``Kitāb al-Shakl al-qattā''} (\emph{Trattato sui quadrilateri}), fu scritto da Nasir al-Din al-Tusi intorno al 1250.}
% https://en.wikipedia.org/wiki/Nasir_al-Din_al-Tusi#Mathematics
\index[persons]{al Din al-Tusi, Nasir}%

%